% Chapter 2: Literature Review

\chapter{Literature Review}

\section{Foundation Models for Earth Observation}

Recent advances in self-supervised learning have produced foundation models
for satellite imagery \citep{jakubik2023foundation, satmae2024}. Prithvi
(IBM-NASA), AlphaEarth (Google), and DINOv2 (Meta) have demonstrated
state-of-the-art performance on land cover classification, change detection,
and semantic segmentation tasks. However, most evaluation has been on
high-resource regions (Europe, US, China); performance in Latin America,
particularly Paraguay, remains under-explored.

\section{Carbon Credit Verification}

Verra VCS \citep{verra2024} and Gold Standard \citep{goldstandard2024} are the
world's largest voluntary carbon markets. Recent work \citep{lamahewage2026alphaearth}
demonstrates that foundation models can estimate forest biomass with
$R^2 > 0.82$ matching ground-based estimates. However, automated verification
remains challenging for projects with sparse ground truth, particularly
smallholder projects in developing countries.

\section{Indigenous Data Governance}

The CARE Principles for Indigenous Data Governance \citep{carroll2020care}
provide a framework for ethical use of indigenous community data. CARE
(Collective Benefit, Authority to Control, Responsibility, Ethics)
complements the FAIR principles and emphasizes community ownership.

\section{Low-Resource Language NLP}

Multilingual NLP for low-resource languages is an active research area
\citep{aguerotorales2023jopara, chiruzzo2023guaspa, kellert2025spanish}.
Spanish-Guaraní code-switching (jopara) is a unique phenomenon that has
received attention \citep{mozillacv2024} but remains challenging for
modern LLMs.

\section{Air Quality Forecasting}

Deep learning approaches \citep{openaq2024} have improved PM2.5 forecasting.
Sentinel-5P atmospheric data \citep{esacopernicus2024} provides global
observations of NO$_2$, SO$_2$, CO, O$_3$, and CH$_4$.

\section{Deforestation Monitoring}

Multi-temporal satellite analysis has become the standard for deforestation
detection \citep{hansen2013high}. MapBiomas Paraguay
\citep{mapbiomas2024paraguay} provides annual land cover maps. NASA FIRMS
\citep{nasafirms2024} provides real-time fire alerts.

\section{Gran Chaco Conservation}

The Gran Chaco is the largest dry forest in South America
\citep{rivas2026gran_chaco}. Defensores del Chaco National Park
\citep{defensores2024} protects key biodiversity. Chagas disease
\citep{who2022chagas, sandon2025chagas} remains a public health concern
linked to deforestation and land-use change.

\section{Latin American Thesis Research}

Paraguay-specific Earth observation research is sparse compared to
neighboring countries (Brazil, Argentina). Notable exceptions include
the work of Cristaldo et al. \citep{cristaldo2024paraguay} on Paraguay
cartography and Agüero-Torales et al. \citep{aguerotorales2023jopara} on
Jopara NLP. This thesis addresses this gap by providing the first
comprehensive multi-paper analysis using foundation models for Paraguay.
